\documentclass[11pt]{article}
\usepackage[letterpaper,margin=0.82in]{geometry}
\usepackage{amsmath,amssymb,amsthm}
\usepackage{booktabs,array,enumitem,microtype,url}
\usepackage[hidelinks]{hyperref}

\newtheorem{theorem}{Theorem}
\newtheorem{proposition}{Proposition}
\newtheorem{lemma}{Lemma}
\newcommand{\Span}{\operatorname{Span}}
\newcommand{\GL}{\operatorname{GL}}
\newcommand{\CO}{\operatorname{CO}}

\title{V22 P4--T02 Source-Local Projective-Conormal Reduction\\
Robustness and Natural-Selection Refuter Stress}
\author{Physicist--Mathematician--Philosopher Parent Synthesis}
\date{15 August 2026}

\begin{document}
\maketitle

\begin{abstract}
The exclusive Refuter class is \texttt{scoped\_obstruction}.  The fixed
source-local projective-conormal reduction candidate remains a mathematically
valid proposal-only construction on its declared coherent submersion-atlas
domain, and the narrow RT006 verdict \texttt{source\_pure\_as\_written} is not
revoked.  The unchanged candidate nevertheless cannot serve as a standalone
robust, naturally selected P4--T02 bridge.  Its integrable domain has no open
ambient $C^1$ protection at the exact fixture; its output-relative arrows
make a valid groupoid but supply no independent admissibility law; all
constant projective lines on $\mathbb R^4$ lie in one proposal-isomorphism
orbit; projectivization removes a $\mathbb Z_2$ coorientation choice; and the
line stabilizer admits no equivariant Lorentz-conformal lift.  These results
justify one candidate-local freeze, not a global no-go.  Seven inherited
freezes and all fourteen Distance-to-GR rows remain unchanged.  One separate
Theoretical Continuation Selector packet is selected and not executed.
\end{abstract}

\section{Fixed candidate and exact scope}

The frozen candidate consists of a smooth four-manifold $M$, a locally finite
cover $\{U_i\}$, submersions $q_i:U_i\to\mathbb R$, and one-dimensional local
diffeomorphisms $\gamma_{ij}$ with
\[
 q_i=\gamma_{ij}\circ q_j.
\]
The lines $\Span(dq_i)$ descend to an unoriented line subbundle
$L_q\subset T^*M$.  Frames adapted to $L_q$ form the principal reduction
$P_H(q)$ for the covector-line stabilizer $H$, and
$\mathrm{Bridge}_{\rm PCR}$ returns $(L_q,P_H(q))$.

This stress holds the complete RT005 candidate and RT006 audit fixed.  It does
not reconstruct or repair the candidate.  It asks whether the unchanged
proposal supplies, from the current bare source reduct, an ambiently robust
domain, an independently admitted arrow law, a quotient-invariant selector,
a coorientation, or a Lorentz-conformal completion.  D7 is not reevaluated,
B2 stays inactive, and P4--T03 remains locked.

\section{Conditional mathematics preserved}

The chain rule gives
\[
 dq_i=\gamma'_{ij}(q_j)dq_j,
\]
with nonzero multiplier, so projective descent, restriction, refinement, and
gluing remain correct on the declared atlas domain.  The line stabilizer in
an adapted basis has first row $(a,0,0,0)$ with $a\ne0$.  Hence
\[
 \dim\GL(4)=16,\qquad \dim H=13,
 \qquad \dim\mathbb P((\mathbb R^4)^*)=3,
\]
and the adapted frames form the claimed proper reduction.

Likewise, proposal arrows defined by
$f^*L_{\rm target}=L_{\rm source}$ close under identities, inverses, and
composition.  This proves a genuine mathematical groupoid.  It does not prove
that the current source ontology independently selects that groupoid as
$\mathrm{EqSrc}$ or an admissibility law.

\section{Exact stress payloads}

\subsection{The coherent-atlas domain is not ambiently protected}

\begin{proposition}[Quantitative non-Frobenius perturbation]
On a fixed compact chart of $\mathbb R^4$, every $C^1$ neighborhood of
$L_0=\Span(dx^0)$ contains both an integrable exact-label line and a
nonintegrable line.  Consequently the candidate's exact-atlas domain is not
open in the ambient $C^1$ space of nonzero line fields at $L_0$.
\end{proposition}

\begin{proof}
For $\epsilon\ne0$, set
\[
 \alpha_\epsilon=dx^0+\epsilon x^1dx^2.
\]
Then
\[
 \alpha_\epsilon\wedge d\alpha_\epsilon
 =\epsilon\,dx^0\wedge dx^1\wedge dx^2\ne0,
\]
so $\ker\alpha_\epsilon$ is not Frobenius integrable.  On
$|x^1|\le R$, the coefficient and first-derivative deviations are bounded by
$|\epsilon|R$ and $|\epsilon|$, respectively, and therefore tend to zero.
In the same neighborhood,
\[
 \beta_\epsilon=d(x^0+\epsilon x^1x^2)
 =dx^0+\epsilon x^2dx^1+\epsilon x^1dx^2
\]
is exact and hence integrable.  Thus no ambient $C^1$ ball about $L_0$ is
contained in the candidate domain.  The RT005 theorem inside coherent exact
labels remains correct; only a universal robustness overread fails.
\end{proof}

Frobenius integrability is diffeomorphism invariant.  The nearby
$\alpha_\epsilon$ proposals therefore cannot be repaired merely by changing
coordinates within the candidate groupoid.  An admissible-variation law that
excludes them would be additional source structure and is not selected by the
bridge output.

\subsection{The proposal quotient removes the strict fixture's separation}

\begin{theorem}[Constant-line proposal-orbit collapse]
All constant projective covector lines on $\mathbb R^4$ lie in one orbit under
orientation-preserving linear source diffeomorphisms.  In particular, the
strict identity-over-$M$ outputs $\Span(dx^0)$ and
$\Span(dx^0+dx^1)$ are unequal but proposal-isomorphic.
\end{theorem}

\begin{proof}
$\GL^+(4,\mathbb R)$ acts transitively on nonzero covectors modulo nonzero
scale: extend any chosen source and target covectors to dual bases, map one
basis to the other, and if necessary reverse two complementary basis vectors
to make the determinant positive without changing the target line.  For the
displayed fixture, the determinant-one shear with first row $(1,1,0,0)$ pulls
$dx^0$ back to $dx^0+dx^1$.
\end{proof}

Thus the RT005 same-$K_\star$ comparison proves only strict embedded
nonfactorization over $\mathrm{id}_M$.  It does not provide a
proposal-isomorphism-invariant discriminator, an occurrence selector, or
physical inequivalence.  This theorem is scoped to constant-line proposals;
it does not claim that every varying foliation presentation is isomorphic.

\subsection{Output-relative arrows are correct but nonselecting}

\begin{lemma}[Closure--admissibility separation]
The line-preserving proposal arrows form a groupoid, while their naturality
has zero independent force for selecting which line, arrow class, or proposal
occurs.
\end{lemma}

\begin{proof}
If $f^*L_2=L_1$ and $g^*L_3=L_2$, then
$(g\circ f)^*L_3=L_1$; identities and inverses are immediate.  But the
naturality equation for the $L$ component is exactly the defining arrow
predicate.  A full-source diffeomorphism that moves $L_1$ is excluded only
after $L_1$ has been supplied.  Citing the resulting naturality cannot derive
the supplied output or elevate the proposal groupoid to adopted
$\mathrm{EqSrc}$.
\end{proof}

\subsection{Projectivization loses coorientation}

\begin{lemma}[Index-two coorientation debt]
Let $H$ stabilize $\Span(dx^0)$ and write the first-row multiplier as
$a\ne0$.  The sign map $H\to\{\pm1\}$ has kernel
$H^+=\{a>0\}$, so selecting a positive conormal ray is an additional
index-two reduction not contained in $P_H(q)$.
\end{lemma}

\begin{proof}
The sign of the product multiplier is the product of signs, giving a
surjective homomorphism.  Its kernel is exactly $H^+$.  Explicitly,
$\operatorname{diag}(-1,-1,1,1)$ has determinant $+1$, preserves the
projective line, and reverses the representative ray.  Hence neither source
orientation nor the projective line chooses a coorientation or physical time
orientation.
\end{proof}

\subsection{A line does not determine conformal geometry}

\begin{theorem}[Equivariant lift obstruction and completion family]
There is no $\GL(4)$-equivariant map from projective covector lines to
Lorentzian conformal classes.  Moreover one fixed line is compatible with an
infinite family of pairwise nonconformal Lorentzian forms.
\end{theorem}

\begin{proof}
A $G$-map $G/H\to G/K$ requires $H$ to lie in a conjugate of $K$.
Here $\dim H=13$, while $\dim\CO(1,3)=7$, so such an inclusion is impossible.
For the pointwise underdetermination, every
\[
 g_{a,b,c}=\operatorname{diag}(-1,a,b,c),\qquad a,b,c>0,
\]
is Lorentzian and makes $dx^0$ timelike.  Normalizing the first coefficient to
$-1$ fixes the conformal scale, so distinct triples $(a,b,c)$ give distinct
conformal classes.  The supplied line therefore determines neither conformal
shape nor physical scale.
\end{proof}

\section{Ten-dimension stress matrix}

\begin{center}
\small
\begin{tabular}{p{0.20\linewidth}p{0.48\linewidth}p{0.22\linewidth}}
\toprule
Dimension & Exact finding & Disposition \\
\midrule
Integrability domain & Non-Frobenius lines occur arbitrarily $C^1$ close to an exact conormal. & ambient robustness obstructed \\
Arrow closure & Output-preserving arrows close as a groupoid. & conditional theorem preserved \\
Arrow force & Naturality repeats the output-preserving premise. & no independent admissibility \\
Proposal quotient & Constant lines form one $\GL^+(4)$ orbit. & strict fixture not quotient-separating \\
Natural selection & Bare-source isotropy has no fixed projective line. & current-source selector absent \\
Coorientation & $H^+\subset H$ has index two. & positive ray and time orientation absent \\
Same reduct & Extra line data changes strict output over identity. & provenance, occurrence still absent \\
Conformal lift & $13$-dimensional $H$ cannot enter $7$-dimensional $\CO(1,3)$. & no equivariant lift \\
Conformal completion & $g_{a,b,c}$ gives infinitely many completions. & line alone underdetermines geometry \\
Authority & Correct proposal math has no empirical, physical, P4, adoption, or Gate force. & all downstream locks fixed \\
\bottomrule
\end{tabular}
\end{center}

\section{Exclusive result and obstruction record}

The exclusive Refuter class is \texttt{scoped\_obstruction}.  The failed
premise is that the unchanged candidate, together with the current bare
source/$K_\star$ assumptions, supplies a source-derived rule that (i) selects
an occurring integrable projective conormal, (ii) protects an ambient
admissible-variation class, (iii) selects a coorientation, and (iv) provides
the missing Lorentz-conformal data.  It supplies none of these four items.

This does not falsify descent, the $H$-reduction, restriction, refinement,
gluing, or groupoid closure on the declared proposal domain.  It does not
revoke \texttt{source\_pure\_as\_written}.  It establishes no global no-go
and no impossibility of a materially enriched future source law.

The candidate-local freeze is
\begin{center}
\small
\path{NDCL-V22-P4T02-B2-PROJECTIVE-CONORMAL-ROBUST-SELECTION-CONFORMAL-LIFT}.
\end{center}
It freezes replay of the unchanged projective-conormal reduction as a
standalone robust, naturally selected, cooriented, Lorentz-conformal P4--T02
bridge.  Its conditional mathematical and diagnostic use remains open.  All
seven inherited freezes remain active and untouched, making eight active
candidate-local freezes after this result.

\section{Distance-to-GR and continuation}

Every expanded Distance-to-GR row remains \texttt{no\_delta}: source ontology,
$\mathrm{EqSrc}$, $\mathrm{RetainH}$, $\mathrm{GenH}$,
$\mathrm{ObsLoc}_{lc}$, $\mathrm{Resp}_{lc}$, $M_{\rm src}$, $g_{\rm eff}$,
matter coupling, Einstein equations, finite-variation robustness, benchmark
promotion, Gate Chair status, and route-freeze status.  The new freeze is a
negative control record, not derivational progress.

Exactly one future packet is selected:
\begin{center}
\small
\path{PKT-V22-P4T02-B2-POST-PROJECTIVE-CONORMAL-REFUTER-THEORETICAL-CONTINUATION-SELECTION-V1}.
\end{center}
Its type and role are
\begin{center}
\small
\path{theoretical_continuation_selector},\qquad
\path{theoretical-continuation-selector@0.1.0}.
\end{center}
Its status is \texttt{selected\_not\_executed}.  It must choose one materially distinct
source-side route or a protected human-gated stop without replaying this or
any inherited freeze.

No source law or ontology object is adopted or rejected.  No physical time,
causal cone, response semantics, conformal geometry, $g_{\rm eff}$, coupling,
Einstein equation, benchmark result, proof authority, publication, push, or
external action is authorized.

\end{document}
